% ============================================================
% BÖLÜM 10: UYGULAMA SÜRECİ
% ============================================================

\chapter{UYGULAMA SÜRECİ}

\section{Geliştirme Aşamaları}

NIDA sisteminin geliştirilmesi, Agile metodolojisi çerçevesinde iteratif bir süreç olarak yürütülmüştür. Geliştirme aşamaları aşağıda detaylandırılmıştır:

\subsection{Aşama 1: Planlama ve Tasarım}

\begin{itemize}
  \item \textbf{Gereksinim Analizi:} Okul psikolojik danışmanlarının ihtiyaçları belirlendi
  \item \textbf{Akademik Literatür Taraması:} HTP temel kaynakları derlendi
  \item \textbf{Sistem Mimarisi Tasarımı:} Katmanlı mimari ve teknoloji seçimleri yapıldı
  \item \textbf{Bilgi Tabanı Şeması:} JSON yapısı ve kural kategorileri tanımlandı
  \item \textbf{Etik Çerçeve:} Kullanım kuralları ve sınırlamalar belirlendi
  \item \textbf{UX/UI Tasarımı:} Wireframe ve prototip çalışmaları yapıldı
\end{itemize}

\subsection{Aşama 2: Bilgi Tabanı Oluşturma}

\begin{itemize}
  \item \textbf{Kaynak Derleme:}
  \begin{itemize}
    \item Buck (1948) - HTP temel kuralları
    \item Hammer (1958) - Klinik uygulama
    \item Machover (1949) - İnsan figürü
    \item Koch (1952) - Ağaç testi
    \item Jolles (1971) - Yorumlama kataloğu
    \item Koppitz (1968, 1984) - Gelişimsel göstergeler
    \item Bolander (1977) - Ağaç analizi
  \end{itemize}
  \item \textbf{Kural Yapılandırma:} 200+ yorumlama kuralı sistematize edildi
  \item \textbf{Çelişki Kaydı:} Çelişkili akademik görüşler dokümante edildi
  \item \textbf{Güven Seviyeleri:} Her kural için güven düzeyi belirlendi
  \item \textbf{Yaş Aralıkları:} Lowenfeld evreleri ile uyumlu kategorizasyon
\end{itemize}

\subsection{Aşama 3: Teknik Geliştirme}

\begin{table}[H]
\centering
\caption{Teknik Geliştirme Adımları}
\label{tab:teknikdev}
\begin{tabular}{p{4cm}p{9cm}}
\toprule
\textbf{Bileşen} & \textbf{Geliştirme Detayları} \\
\midrule
Proje Kurulumu & Next.js 16 + TypeScript + Tailwind CSS \\
Firebase Entegrasyonu & Authentication, Firestore, Storage kurulumu \\
Kullanıcı Arayüzü & Login, Register, Dashboard, Analyze sayfaları \\
API Geliştirme & /api/analyze endpoint, Gemini entegrasyonu \\
Bilgi Tabanı & knowledge-base.ts modülü, getKBRules() fonksiyonu \\
Prompt Mühendisliği & Türkçe/İngilizce çift dilli prompt optimizasyonu \\
Raporlama & PDF export fonksiyonalitesi \\
Deployment & Vercel deployment, environment variables \\
\bottomrule
\end{tabular}
\end{table}

\subsection{Aşama 4: Test ve İyileştirme}

\begin{itemize}
  \item \textbf{Fonksiyonel Testler:} Tüm kullanım senaryoları test edildi
  \item \textbf{Dil Testleri:} Türkçe ve İngilizce çıktılar doğrulandı
  \item \textbf{Performans Testi:} API yanıt süreleri ölçüldü (ortalama 7-8 saniye)
  \item \textbf{Güvenlik Testi:} Kimlik doğrulama ve yetkilendirme kontrolleri
  \item \textbf{Kullanılabilirlik:} Responsive tasarım ve erişilebilirlik kontrolü
  \item \textbf{Hata Giderme:} Tespit edilen hatalar düzeltildi
\end{itemize}

\section{Teknoloji Seçimi Gerekçeleri}

\begin{table}[H]
\centering
\caption{Teknoloji Seçimleri ve Gerekçeleri}
\label{tab:teknolojisecim}
\begin{tabular}{p{3cm}p{10cm}}
\toprule
\textbf{Teknoloji} & \textbf{Seçim Gerekçesi} \\
\midrule
Next.js 16 & Modern React framework, App Router, Server Components, kolay deployment \\
TypeScript & Tip güvenliği, kod kalitesi, IDE desteği, bakım kolaylığı \\
Tailwind CSS & Hızlı prototipleme, tutarlı tasarım sistemi, düşük dosya boyutu \\
Firebase & Hızlı başlangıç, gerçek zamanlı veritabanı, entegre auth, ücretsiz tier \\
Gemini 2.0 Flash & Çok modlu analiz, Türkçe dil desteği, maliyet etkinliği, düşük gecikme \\
Vercel & Next.js için optimize edilmiş hosting, otomatik ölçekleme, kolay CI/CD \\
\bottomrule
\end{tabular}
\end{table}

\section{Prompt Mühendisliği}

Sistemin başarısı, büyük ölçüde yapay zekâ modeline gönderilen prompt'un kalitesine bağlıdır. NIDA'da kullanılan prompt yapısı şu bileşenlerden oluşmaktadır:

\subsection{Prompt Yapısı}

\begin{enumerate}
  \item \textbf{Rol Tanımı:} Modelin HTP uzmanı olarak konumlandırılması
  \begin{quote}
  ``Sen projektif çizim analizi, özellikle House-Tree-Person (HTP) testi konusunda uzmanlaşmış bir çocuk psikoloğusun.''
  \end{quote}

  \item \textbf{Bilgi Tabanı Enjeksiyonu:} İlgili yorumlama kurallarının prompt'a dahil edilmesi

  \item \textbf{Etik Hatırlatmalar:}
  \begin{itemize}
    \item Kesin tanı ifadeleri kullanılmaması
    \item Yumuşatıcı dil kullanımı (``gösterebilir'', ``işaret edebilir'')
    \item Her yorum için güven seviyesi belirtilmesi
    \item Kaynak zorunluluğu
  \end{itemize}

  \item \textbf{Çıktı Formatı:} Yapılandırılmış JSON şeması

  \item \textbf{Dil Ayarı:} Türkçe veya İngilizce çıktı kontrolü
\end{enumerate}

\subsection{Prompt Optimizasyonu}

Prompt geliştirme sürecinde aşağıdaki iterasyonlar gerçekleştirildi:

\begin{enumerate}
  \item \textbf{V1:} Temel İngilizce prompt - Dil karışıklığı sorunu
  \item \textbf{V2:} Türkçe talimatlar eklendi - Kısmi iyileşme
  \item \textbf{V3:} Tam Türkçe prompt yazıldı - İstenilen sonuç
  \item \textbf{V4:} Güven seviyesi ve kaynak zorunluluğu eklendi
  \item \textbf{V5:} JSON çıktı formatı optimize edildi
\end{enumerate}

\section{Karşılaşılan Zorluklar ve Çözümler}

\begin{table}[H]
\centering
\caption{Geliştirme Sürecinde Karşılaşılan Zorluklar}
\label{tab:zorluklar}
\begin{tabular}{p{4cm}p{4cm}p{5cm}}
\toprule
\textbf{Zorluk} & \textbf{Etki} & \textbf{Çözüm} \\
\midrule
Gemini API Kota Limiti & 429 hatası, hizmet kesintisi & Ücretli plan aktifleştirildi \\
Türkçe Çıktı Sorunu & İngilizce yanıtlar & Prompt tamamen Türkçe yazıldı \\
Firebase Index Hatası & Sorgu başarısız & Composite index oluşturuldu \\
TypeScript Tip Hataları & Build başarısız & Tip tanımları güncellendi \\
Build Sırasında Env & Firebase init hatası & Null check eklendi \\
\bottomrule
\end{tabular}
\end{table}

\section{Deployment ve Canlıya Alma}

\subsection{Deployment Süreci}

\begin{enumerate}
  \item GitHub repository oluşturuldu
  \item Vercel ile GitHub entegrasyonu sağlandı
  \item Environment variables tanımlandı
  \item Otomatik deployment aktifleştirildi
  \item Custom domain yapılandırması (opsiyonel)
\end{enumerate}

\subsection{Canlı Ortam Bilgileri}

\begin{itemize}
  \item \textbf{Vercel URL:} https://app-baristiran.vercel.app
  \item \textbf{GitHub:} https://github.com/tialkan/nida
  \item \textbf{API Endpoint:} /api/analyze
  \item \textbf{Desteklenen Diller:} Türkçe (varsayılan), İngilizce
\end{itemize}

\newpage
